\section{Configuration Phase}

The configuration phase runs before inference and is responsible for converting
the serialized parameter stream into aligned per-layer writes. In this design,
that work is intentionally split into two levels. First, the top-level
\texttt{config\_manager} reshapes the incoming byte stream and separates weight
payloads from threshold payloads. Then each \texttt{bnn\_layer} uses a local
\texttt{config\_controller} path to distribute those writes across replicated
RAM banks.

\subsection{Configuration Manager}

The top-level configuration datapath is implemented in
\texttt{rtl/config\_manager.sv} with support from
\texttt{tkeep\_byte\_compactor}, \texttt{vw\_buffer},
\texttt{config\_manager\_parser}, and \texttt{config\_manager\_pad\_fsm}. Its
job is to accept the external configuration stream, compact away invalid bytes,
recover message structure, and emit write-oriented streams that are already
separated into the weight and threshold paths.

\begin{strip}
  \centering
  \resizebox{!}{0.41\textheight}{\begin{tikzpicture}[
  >=Latex,
  block/.style={
    draw,
    thick,
    rounded corners=1pt,
    minimum width=60mm,
    minimum height=17mm,
    align=center,
    font=\normalsize
  },
  wideblock/.style={
    block,
    minimum width=74mm
  },
  data/.style={->, thick},
  bus/.style={
    ->,
    thick,
    postaction={decorate},
    decoration={
      markings,
      mark=at position 0.55 with {
        \draw[thick] (-2.2pt,3pt) -- (2.2pt,-3pt);
      }
    }
  },
  ctrl/.style={->, thick, dashed},
  iface/.style={font=\large\bfseries\ttfamily, fill=white, inner sep=1pt},
  sig/.style={font=\normalsize\ttfamily, fill=white, inner sep=1pt},
  lab/.style={font=\normalsize, fill=white, inner sep=1pt}
]

\node[font=\Large\bfseries] at (0,3.3) {Config Manager};

\node[block] (compact) at (0,0)
  {TKEEP Byte Compactor\\[-1pt]\texttt{tkeep\_byte\_compactor}};

\node[block] (vw) at (0,-4.1)
  {VW Buffer\\[-1pt]\texttt{vw\_buffer}};

\node[block] (cfgfifo) at (0,-8.2)
  {Config Byte FIFO\\[-1pt]\texttt{fifo\_config\_bytes}};

\node[wideblock] (parser) at (0,-12.4)
  {Parser\\[-1pt]\texttt{config\_manager\_parser}};

\node[block] (wfifo) at (-11.0,-18.3)
  {Weight-Byte FIFO\\[-1pt]\texttt{fifo\_weights\_bytes}};

\node[block] (pad) at (0,-18.3)
  {Pad FSM\\[-1pt]\texttt{config\_manager\_pad\_fsm}};

\node[block] (tfifo) at (11.0,-18.3)
  {Threshold FIFO\\[-1pt]\texttt{fifo\_thresholds}};

\node[block] (pack) at (0,-23.2)
  {Layer Packers\\[-1pt]\texttt{fifo\_packer[i]}};

% Top-level config inputs
\coordinate (cfgvTop) at ([xshift=-50mm,yshift=18mm]compact.north);
\coordinate (cfgvTurnA) at ([xshift=-50mm,yshift=8mm]compact.north);
\coordinate (cfgvTurnB) at ([xshift=-20mm,yshift=8mm]compact.north);
\draw[data] (cfgvTop) -- (cfgvTurnA) -- (cfgvTurnB) -- ([xshift=-20mm]compact.north);
\node[iface] at ($(cfgvTop)+(0,2.5mm)$) {\texttt{config\_valid}};

\draw[bus] ([yshift=18mm]compact.north) -- (compact.north);
\node[iface] at ([yshift=20.5mm]compact.north) {\texttt{config\_data}};

\coordinate (cfgkTop) at ([xshift=50mm,yshift=18mm]compact.north);
\coordinate (cfgkTurnA) at ([xshift=50mm,yshift=8mm]compact.north);
\coordinate (cfgkTurnB) at ([xshift=20mm,yshift=8mm]compact.north);
\draw[bus] (cfgkTop) -- (cfgkTurnA) -- (cfgkTurnB) -- ([xshift=20mm]compact.north);
\node[iface] at ($(cfgkTop)+(0,2.5mm)$) {\texttt{config\_keep}};

% Compactor -> VW buffer
\draw[data] ([xshift=-22mm]compact.south) -- ([xshift=-22mm]vw.north)
  node[pos=0.64,left,sig] {\texttt{compact\_wr\_en}};

\draw[bus] (compact.south) -- (vw.north)
  node[pos=0.38,right,sig,xshift=2pt] {\texttt{compact\_wr\_data}};

\draw[bus] ([xshift=22mm]compact.south) -- ([xshift=22mm]vw.north)
  node[pos=0.64,right,sig] {\texttt{compact\_total\_bytes}};

% VW buffer -> config byte FIFO
\draw[data] ([xshift=-15mm]vw.south) -- ([xshift=-15mm]cfgfifo.north)
  node[pos=0.50,left,sig] {\texttt{cfg\_vw\_rd\_en}};

\draw[bus] ([xshift=15mm]vw.south) -- ([xshift=15mm]cfgfifo.north)
  node[pos=0.50,right,sig] {\texttt{cfg\_vw\_rd\_data}};

% Config byte FIFO -> parser
\draw[bus] (cfgfifo.south) -- (parser.north)
  node[pos=0.58,right,sig] {\texttt{cfg\_byte\_data}};

% config_ready routed back up as interface signal
\coordinate (cfgreadyA) at ($(cfgfifo.east)+(5.0,0)$);
\coordinate (cfgreadyTop) at ($(cfgreadyA |- compact.north)+(0,18mm)$);

\draw[thick] (cfgfifo.east) -- (cfgreadyA);
\draw[data] (cfgreadyA) -- (cfgreadyTop);
\node[iface] at ($(cfgreadyTop)+(0,2.3mm)$) {\texttt{config\_ready}};

% Parser -> payload destinations
\draw[bus] ([xshift=-23mm]parser.south) -- (wfifo.north)
  node[pos=0.48,above left,lab] {weight payload bytes};

\draw[bus] ([xshift=23mm]parser.south) -- (tfifo.north)
  node[pos=0.48,above right,lab] {threshold payload bytes};

% Parser -> pad FSM control
\draw[data] ([xshift=-20mm]parser.south) -- ([xshift=-20mm]pad.north)
  node[pos=0.42,left,sig] {\texttt{msg\_type}};

\draw[bus] (parser.south) -- (pad.north)
  node[pos=0.42,right,sig,xshift=2pt] {\texttt{layer\_id}};

\draw[data] ([xshift=20mm]parser.south) -- ([xshift=20mm]pad.north)
  node[pos=0.42,right,sig] {\texttt{payload\_start}};

% Pad FSM -> parser stall
\coordinate (stallA) at ($(pad.north east)+(1.5,0)$);
\coordinate (stallB) at ($(stallA)+(0,4.8)$);
\coordinate (stallC) at ($(parser.east)+(1.2,0)$);

\draw[ctrl] (pad.north east) -- (stallA) -- (stallB) -- (stallC) -- (parser.east);
\node[sig,right] at ($(stallA)!0.55!(stallB)$) {\texttt{stall}};

% Weight path
\draw[bus] (wfifo.east) -- (pad.west)
  node[pos=0.52,above,sig] {\texttt{w\_byte\_data}};

\draw[bus] (pad.south) -- (pack.north)
  node[pos=0.55,right,sig] {\texttt{data}};

% Threshold outputs
\coordinate (thDataA) at ([xshift=-12mm]tfifo.south);
\coordinate (thDataB) at ($(thDataA)+(0,-1.0)$);
\coordinate (thDataC) at ($(thDataB)+(-18mm,0)$);
\draw[bus] (thDataA) -- (thDataB) -- (thDataC) -- ++(0,-1.6);
\node[below,iface] at ($(thDataC)+(0,-1.6)$) {\texttt{threshold\_wr\_data}};

\coordinate (thEnA) at ([xshift=12mm]tfifo.south);
\coordinate (thEnB) at ($(thEnA)+(0,-1.0)$);
\coordinate (thEnC) at ($(thEnB)+(18mm,0)$);
\draw[data] (thEnA) -- (thEnB) -- (thEnC) -- ++(0,-1.6);
\node[below,iface] at ($(thEnC)+(0,-1.6)$) {\texttt{threshold\_wr\_en}};

% Weight outputs
\draw[bus] ([xshift=-17mm]pack.south) -- ++(0,-2.6)
  node[below,iface] {\texttt{weight\_wr\_data}};

\draw[data] ([xshift=17mm]pack.south) -- ++(0,-2.6)
  node[below,iface] {\texttt{weight\_wr\_en}};

\end{tikzpicture}
}
  \captionof{figure}{Block diagram of the \texttt{config\_manager} datapath and
  its primary interface signals.}
  \Description{Vertical block diagram of the configuration manager showing
  byte compaction, repacking, parsing, weight padding, threshold packing, and
  the external write interfaces that leave the module.}
  \label{fig:config-manager}
\end{strip}

The front end of this path is byte-oriented because the external interface can
arrive with sparse valid-byte patterns. The parser then restores the higher
level meaning of the stream by decoding the message header, choosing the weight
or threshold destination, and generating the control metadata needed by the
padding logic. On the weight side, the pad state machine inserts synthetic
bytes so the downstream layer packers can emit aligned write words.

\clearpage

\subsection{Layer-Local Write Distribution}

After the top-level parser finishes that global reshaping step, each
\texttt{bnn\_layer} still has to fan the write stream out to replicated local
memories. That is the purpose of \texttt{config\_controller} and the
surrounding staging registers in \texttt{rtl/bnn\_layer.sv}. The controller
generates the bank-select and address signals, while the aligned data stream
bypasses the controller datapath and is retimed for fanout.

\begin{figure*}[t]
  \centering
  \resizebox{!}{0.31\textheight}{\begin{tikzpicture}[
  >=Latex,
  block/.style={
    draw,
    thick,
    rounded corners=1pt,
    minimum width=58mm,
    minimum height=18mm,
    align=center,
    font=\normalsize
  },
  wideblock/.style={
    block,
    minimum width=72mm
  },
  data/.style={->, thick},
  bus/.style={
    ->,
    thick,
    postaction={decorate},
    decoration={
      markings,
      mark=at position 0.55 with {
        \draw[thick] (-2.2pt,3pt) -- (2.2pt,-3pt);
      }
    }
  },
  iface/.style={font=\large\bfseries\ttfamily, fill=white, inner sep=1pt},
  sig/.style={font=\normalsize\ttfamily, fill=white, inner sep=1pt},
  lab/.style={font=\normalsize, fill=white, inner sep=1pt}
]

\node[font=\Large\bfseries] at (0,3.4) {Layer Config Write Distribution};

\node[block] (wstage) at (-8.2,0)
  {Weight Input Regs\\[-1pt]\textit{ingress timing cut}};

\node[block] (tstage) at (8.2,0)
  {Threshold Input Regs\\[-1pt]\textit{ingress timing cut}};

\node[wideblock] (ctrlblk) at (0,-5.1)
  {Config Controller\\[-1pt]\texttt{config\_controller}\\[-1pt]\textit{bank select + address generation}};

\node[block] (wfan) at (-8.2,-10.7)
  {Weight Fanout Regs\\[-1pt]\textit{timing cut before RAM banks}};

\node[block] (tfan) at (8.2,-10.7)
  {Threshold Fanout Regs\\[-1pt]\textit{timing cut before RAM banks}};

\node[block] (wram) at (-8.2,-16.0)
  {Weight RAM Banks\\[-1pt]\texttt{u\_w\_ram[gi]}\\[-1pt]\textit{one bank per NP}};

\node[block] (tram) at (8.2,-16.0)
  {Threshold RAM Banks\\[-1pt]\texttt{u\_t\_ram[gi]}\\[-1pt]\textit{one bank per NP}};

% Top-level layer-local inputs
\draw[bus] ([xshift=-12mm,yshift=18mm]wstage.north) -- ([xshift=-12mm]wstage.north);
\node[iface] at ([xshift=-12mm,yshift=20.5mm]wstage.north) {\texttt{w\_data}};

\draw[data] ([xshift=12mm,yshift=18mm]wstage.north) -- ([xshift=12mm]wstage.north);
\node[iface] at ([xshift=12mm,yshift=20.5mm]wstage.north) {\texttt{w\_en}};

\draw[bus] ([xshift=-12mm,yshift=18mm]tstage.north) -- ([xshift=-12mm]tstage.north);
\node[iface] at ([xshift=-12mm,yshift=20.5mm]tstage.north) {\texttt{t\_data}};

\draw[data] ([xshift=12mm,yshift=18mm]tstage.north) -- ([xshift=12mm]tstage.north);
\node[iface] at ([xshift=12mm,yshift=20.5mm]tstage.north) {\texttt{t\_en}};

% Stage -> controller enables, orthogonal
\coordinate (wenA) at ([xshift=14mm]wstage.south);
\coordinate (wenB) at ($(wenA)+(0,-1.2)$);
\coordinate (wenC) at ($(ctrlblk.north west)+(-6mm,1.2)$);
\draw[data] (wenA) -- (wenB) -- (wenC) -- ([xshift=-18mm]ctrlblk.north)
  node[pos=0.56,above,sig] {\texttt{w\_en\_r}};

\coordinate (tenA) at ([xshift=-14mm]tstage.south);
\coordinate (tenB) at ($(tenA)+(0,-1.2)$);
\coordinate (tenC) at ($(ctrlblk.north east)+(6mm,1.2)$);
\draw[data] (tenA) -- (tenB) -- (tenC) -- ([xshift=18mm]ctrlblk.north)
  node[pos=0.56,above,sig] {\texttt{t\_en\_r}};

% Staged write data bypasses controller and goes straight down
\draw[bus] ([xshift=-14mm]wstage.south) -- ([xshift=-14mm]wfan.north)
  node[pos=0.50,left,sig] {\texttt{w\_data\_r}};

\draw[bus] ([xshift=14mm]tstage.south) -- ([xshift=14mm]tfan.north)
  node[pos=0.50,right,sig] {\texttt{t\_data\_r}};

% Controller outputs -> fanout timing cut, orthogonal
\coordinate (wselA) at ([xshift=-21mm]ctrlblk.south);
\coordinate (wselB) at ($(wselA)+(0,-1.2)$);
\coordinate (wselT) at ([xshift=16mm]wfan.north);
\coordinate (wselC) at (wselT |- wselB);
\draw[data] (wselA) -- (wselB) -- (wselC) -- (wselT);
\node[above,sig] at ($(wselB)!0.5!(wselC)$) {\texttt{w\_sel}};

\coordinate (waddrA) at ([xshift=-7mm]ctrlblk.south);
\coordinate (waddrB) at ($(waddrA)+(0,-2.1)$);
\coordinate (waddrT) at (wfan.north);
\coordinate (waddrC) at (waddrT |- waddrB);
\draw[bus] (waddrA) -- (waddrB) -- (waddrC) -- (waddrT);
\node[above,sig] at ($(waddrB)!0.5!(waddrC)$) {\texttt{w\_addr[i]}};

\coordinate (taddrA) at ([xshift=7mm]ctrlblk.south);
\coordinate (taddrB) at ($(taddrA)+(0,-2.1)$);
\coordinate (taddrT) at (tfan.north);
\coordinate (taddrC) at (taddrT |- taddrB);
\draw[bus] (taddrA) -- (taddrB) -- (taddrC) -- (taddrT);
\node[above,sig] at ($(taddrB)!0.5!(taddrC)$) {\texttt{t\_addr[i]}};

\coordinate (tselA) at ([xshift=21mm]ctrlblk.south);
\coordinate (tselB) at ($(tselA)+(0,-1.2)$);
\coordinate (tselT) at ([xshift=-16mm]tfan.north);
\coordinate (tselC) at (tselT |- tselB);
\draw[data] (tselA) -- (tselB) -- (tselC) -- (tselT);
\node[above,sig] at ($(tselB)!0.5!(tselC)$) {\texttt{t\_sel}};

% config_done routed upward as a layer-local control signal
\coordinate (cfgdoneTop) at ($(ctrlblk.north)+(0,6.4)$);
\draw[data] (ctrlblk.north) -- (cfgdoneTop);
\node[iface] at ($(cfgdoneTop)+(0,2.3mm)$) {\texttt{done}};

% Fanout timing cut -> RAM banks
\draw[data] ([xshift=-16mm]wfan.south) -- ([xshift=-16mm]wram.north)
  node[pos=0.52,left,sig] {\texttt{w\_sel\_p[i]}};

\draw[bus] (wfan.south) -- (wram.north)
  node[pos=0.52,right,sig] {\texttt{w\_addr\_p[i]}};

\draw[bus] ([xshift=16mm]wfan.south) -- ([xshift=16mm]wram.north)
  node[pos=0.52,right,sig] {\texttt{w\_data\_p}};

\draw[data] ([xshift=-16mm]tfan.south) -- ([xshift=-16mm]tram.north)
  node[pos=0.52,left,sig] {\texttt{t\_sel\_p[i]}};

\draw[bus] (tfan.south) -- (tram.north)
  node[pos=0.52,right,sig] {\texttt{t\_addr\_p[i]}};

\draw[bus] ([xshift=16mm]tfan.south) -- ([xshift=16mm]tram.north)
  node[pos=0.52,right,sig] {\texttt{t\_data\_p}};

\end{tikzpicture}
}
  \caption{Layer-local config-write distribution around
  \texttt{config\_controller}. The highlighted fanout pipeline registers are
  the timing-cut stage that drives the per-bank weight and threshold RAMs.}
  \Description{Block diagram of the layer-local configuration write path
  showing staged weight and threshold inputs, the config controller, timing-cut
  fanout registers, and replicated RAM banks for weight and threshold storage.}
  \label{fig:config-controller-block}
\end{figure*}

Those fanout pipeline registers are a central part of the implementation, not a
minor detail. They cut the long controller-to-BRAM distribution path into a
manageable stage while preserving a steady stream of writes after the pipeline
fills. That is why this report treats the local write-distribution block as its
own step in the overall configuration phase.

\clearpage
